\documentclass{beamer}
\usepackage{amsmath, amssymb}
\usepackage{bm}

\title{Mathematical Formulation of BERT\\(Masked Language Modeling)}
\author{Qingfeng Liu}
\date{\today }

\begin{document}
	
	\frame{\titlepage}
	
	% --- Notation ---
	\begin{frame}{Notation}
		\begin{itemize}
			\item $T$: Sequence length
			\item $V$: Vocabulary size
			\item $d$: Dimensionality of hidden states
			\item $L$: Number of Transformer encoder layers
			\item $\mathcal{M} \subset \{1, \dots, T\}$: Set of masked token positions
			\item $y_i \in \{0, \dots, V{-}1\}$: Ground-truth token index at position $i \in \mathcal{M}$
			\item $X \in \mathbb{R}^{T \times d}$: Input embeddings with positional encoding
			\item $H^{(\ell)} \in \mathbb{R}^{T \times d}$: Output of layer $\ell$
			\item $A^{(\ell)} \in \mathbb{R}^{T \times d}$: Attention output of layer $\ell$
			\item $W \in \mathbb{R}^{V \times d}$,\quad $b \in \mathbb{R}^V$: Output projection weights and bias
		\end{itemize}
	\end{frame}
	
	% --- What is a Token ---
	\begin{frame}{What is a Token?}
		\begin{itemize}
			\item A token is the smallest text unit handled by BERT.
			\item The vocabulary $\mathcal{V}$ is a fixed-size set (e.g., 30,522 tokens).
			\item Tokens are based on WordPiece subword units.
			\item Example: \texttt{unbelievable} $\rightarrow$ [\texttt{"un"}, \texttt{"\#\#believ"}, \texttt{"\#\#able"}]
		\end{itemize}
	\end{frame}
	
	% --- Token IDs and Embedding ---
	\begin{frame}{Token IDs and Embedding Vectors}
		\begin{itemize}
			\item Each token is mapped to a vocabulary index: $x_t \in \{0, 1, \dots, V{-}1\}$
			\item A learnable embedding matrix $E \in \mathbb{R}^{V \times d}$ assigns:
			\[
			e_t = E[x_t] \in \mathbb{R}^d
			\]
			\item $e_t$ is the vector representation of token $x_t$
		\end{itemize}
	\end{frame}
	
	% --- Embedding Construction ---
	\begin{frame}{Input Embedding Construction}
		Token embedding $e_t$ is combined with positional embedding $p_t$:
		\[
		x_t = e_t + p_t
		\]
		
		All token embeddings:
		\[
		X = E[x_1, x_2, \dots, x_T] + P, \quad X \in \mathbb{R}^{T \times d}
		\]
		\begin{itemize}
			\item $P$: position embedding matrix (fixed or learnable)
		\end{itemize}
	\end{frame}
	
	% --- What is the Embedding Space ---
	\begin{frame}{What is the Embedding Space?}
		\begin{itemize}
			\item The embedding space is the $d$-dimensional real vector space $\mathbb{R}^d$
			\item Each token is represented as a vector $e_t \in \mathbb{R}^d$
			\item Similar tokens are placed closer together in this space
			\item This structure is learned during training based on language context
		\end{itemize}
	\end{frame}
	
	% --- Who Determines the Space? ---
	\begin{frame}{Who Determines the Embedding Space?}
		\begin{itemize}
			\item Humans design the dimensions $d$ and vocabulary size $V$
			\item The matrix $E$ is randomly initialized (e.g., from $\mathcal{N}(0, \sigma^2 I)$)
			\item It is optimized via backpropagation using task losses:
			\[
			E_k \leftarrow E_k - \eta \cdot \nabla_{E_k} \mathcal{L}
			\]
			\item The space is shaped to reflect semantic similarity through learning
		\end{itemize}
	\end{frame}
	
	% --- Summary ---
	\begin{frame}{Summary}
		\begin{itemize}
			\item The embedding space is not predefined — it is learned through training
			\item Tokens acquire meaning via their placement in a distributed representation space
			\item The quality of the embedding space defines BERT's ability to model language
		\end{itemize}
	\end{frame}
	% --- Step 1 ---
	\begin{frame}{Step 1: Input Embedding}
		\[
		X = \text{TokenEmbed}(x_1, \dots, x_T) + \text{PositionEmbed}(1, \dots, T)
		\]
	\end{frame}
	
	% --- Step 2 ---
	\begin{frame}{Step 2: Transformer Encoder ($\ell = 1, \dots, L$)}
		\[
		\begin{aligned}
			Q^{(\ell)} &= H^{(\ell-1)} W^{Q(\ell)} \\
			K^{(\ell)} &= H^{(\ell-1)} W^{K(\ell)} \\
			V^{(\ell)} &= H^{(\ell-1)} W^{V(\ell)} \\
			A^{(\ell)} &= \text{softmax}\left( \frac{Q^{(\ell)} K^{(\ell)\top}}{\sqrt{d}} \right) V^{(\ell)} \\
			\tilde{H}^{(\ell)} &= \text{LayerNorm}(A^{(\ell)} + H^{(\ell-1)}) \\
			H^{(\ell)} &= \text{LayerNorm}(\text{FFN}(\tilde{H}^{(\ell)}) + \tilde{H}^{(\ell)})
		\end{aligned}
		\]
	\end{frame}
	
	% --- FFN ---
	\begin{frame}{Feed Forward Network (FFN)}
		\[
		\text{FFN}(x) = \text{ReLU}(x W_1 + b_1) W_2 + b_2
		\]
		\begin{itemize}
			\item Applied independently to each token
			\item Typically: hidden size $d_{\text{ff}} = 4d$
		\end{itemize}
	\end{frame}
	
	% --- Step 3 ---
	\begin{frame}{Step 3: Final Output}
		\[
		H = H^{(L)} \in \mathbb{R}^{T \times d}
		\]
		\begin{itemize}
			\item Used for prediction at masked positions
		\end{itemize}
	\end{frame}
	
	% --- Step 4 ---
	\begin{frame}{Step 4: Output Layer (Vocabulary Projection)}
		\[
		\forall i \in \mathcal{M},\quad \hat{y}_i = \text{softmax}(W H_i + b) \in \mathbb{R}^V
		\]
		\begin{itemize}
			\item $\hat{y}_i[k]$: Probability of token $k$ at position $i$
		\end{itemize}
	\end{frame}
	
	% --- Step 5 ---
	\begin{frame}{Step 5: Loss Function (Cross Entropy)}
		\[
		\mathcal{L}_{\text{MLM}} = - \sum_{i \in \mathcal{M}} \log \hat{y}_i[y_i]
		\]
	\end{frame}
	
	% --- Final Loss ---
	\begin{frame}{Final Loss Expression}
		\[
		\boxed{
			\mathcal{L}_{\text{BERT}} = - \sum_{i \in \mathcal{M}} \log \left( \text{softmax}(W H_i + b)[y_i] \right)
		}
		\]
	\end{frame}
	
	% --- Terminology ---
	\begin{frame}{Terminology Summary}
		\begin{itemize}
			\item $H^{(\ell)}$: Hidden state at layer $\ell$
			\item $A^{(\ell)}$: Attention output at layer $\ell$
			\item FFN: Feed Forward Network (2-layer MLP with ReLU)
		\end{itemize}
	\end{frame}
	
	
	% --- Basic Equation ---
	\begin{frame}{Basic Formula of Attention}
		\[
		\text{Attention}(Q, K, V) = \mathrm{softmax}\left( \frac{QK^\top}{\sqrt{d_k}} \right) V
		\]
		\begin{itemize}
			\item \( Q \in \mathbb{R}^{T \times d_k} \): Queries
			\item \( K \in \mathbb{R}^{T \times d_k} \): Keys
			\item \( V \in \mathbb{R}^{T \times d_v} \): Values
		\end{itemize}
	\end{frame}
	
	% --- Inner Product Interpretation ---
	\begin{frame}{Step 1: Inner Product Similarity}
		\[
		S = QK^\top \in \mathbb{R}^{T \times T}
		\]
		\begin{itemize}
			\item Each element \( S_{ij} = \langle q_i, k_j \rangle \) is a dot product
			\item Measures similarity between query \( q_i \) and key \( k_j \)
			\item Geometrically: projection of \( q_i \) onto key space
		\end{itemize}
	\end{frame}
	
	% --- Softmax Normalization ---
	\begin{frame}{Step 2: Softmax Normalization}
		\[
		A = \mathrm{softmax}\left( \frac{QK^\top}{\sqrt{d_k}} \right)
		\]
		\begin{itemize}
			\item Scaling by \( \sqrt{d_k} \) controls sharpness of softmax
			\item Each row of \( A \) becomes a probability distribution
			\item Represents attention weights assigned to values
		\end{itemize}
	\end{frame}
	
	% --- Weighted Sum as Expectation ---
	\begin{frame}{Step 3: Weighted Average of Values}
		\[
		Z = AV \in \mathbb{R}^{T \times d_v}
		\]
		\begin{itemize}
			\item Each output vector \( z_i \) is a weighted sum of all value vectors
			\item Weights determined by similarity between \( q_i \) and all \( k_j \)
			\item Can be interpreted as a soft "semantic lookup"
		\end{itemize}
	\end{frame}
	
	% --- Full Interpretation ---
	\begin{frame}{Full Mathematical Interpretation}
		\[
		\boxed{
			\text{Attention}(q_i, K, V) = \sum_{j=1}^T \mathrm{softmax}_j\left( \frac{q_i^\top k_j}{\sqrt{d_k}} \right) \cdot v_j
		}
		\]
		\begin{itemize}
			\item Weighted average of values based on similarity scores
			\item Nonlinear weighted linear transformation
		\end{itemize}
	\end{frame}
	
	% --- Geometric View ---
	\begin{frame}{Geometric Interpretation}
		\begin{itemize}
			\item Query \( q_i \) is compared with all keys \( k_j \)
			\item Similar keys receive higher attention weights
			\item Output vector is pulled toward directions of semantically relevant values
			\item Attention reshapes input space dynamically via similarity structure
		\end{itemize}
	\end{frame}
	
\end{document}